\documentclass[12pt]{article}
\usepackage[utf8]{inputenc}
\usepackage[T1]{fontenc}
\usepackage[serbian]{babel}
\usepackage{indentfirst}
\usepackage{geometry}
\geometry{a4paper, margin=2.5cm}

\begin{document}

\begin{center}
    \textit{Maska je morala biti polovična. Cvet je morao da bude ruža.}\\
    \textit{Nije mogao da postoji kao takav lik bez toga.}\\
    \textit{Ali, da li je on zapravo ikada bio neko osim tih osobina?}\\[0.5cm]

    {\Large \textbf{Erikova uloga Fantoma -- Da li smo ikada zaista shvatili ko je on, ili samo ono što je bio primoran da postane?}} \\[0.4cm]

    {\small \textit{Autor: Dorotea Borenović, diplomirani profesionalni plesač, BŠNS}} \\
    {\small \textit{Predmet: istorija umetnosti}}
\end{center}

\vspace{0.8cm}

Fantom iz opere je tajanstveni lik koji živi ispod pariske Opere, skriven od sveta i svetla. Njegovo pravo ime je Erik, ali skoro niko ga tako ne zove. Za sve ostale, on je samo Fantom. Senka, glas iz tame, biće koje postoji između stvarnosti i mita. On ne živi među ljudima, već ispod njih, u podzemlju opere, u prostoru gde nema publike, aplauza ni prihvatanja. Već tu se vidi da Fantom nije samo lik, već simbol nekoga ko je izbačen sa pozornice života. Neko ko ne želi da se prikazuje. 

Jedan od najprepoznatljivijih delova Fantomovog lika jeste njegova polovična maska. Ona ne skriva celo lice, već samo deo, kao da ni sam Erik nikada nije mogao potpuno da se sakrije, ali ni potpuno da se pokaže. Ta maska nije samo fizička zaštita, već znak njegove unutrašnje podele. On je istovremeno genije i čudovište u očima drugih. Maska mu omogućava da postoji u svetu koji ga ne prihvata, ali ga istovremeno zauvek odvaja od tog sveta. Ideja polovične maske je tek došla u operi i mjuziklu, u knjizi je ona bila cela. Vizuelno ta maska govori mnogo više ako je polovična - pokazuje kakva je osoba zapravo fantom. Mi njega ne možemo posmatrati u potpunosti kao negativca. U njemu postoji dobrota, ljudskost i nežnost i to se vidi kad pokazuje ljubav prema Kristini. Jednostavno se ne može posmatrati kao potpuni negativac.

Ruže koje Fantom ostavlja, posebno one namenjene Kristini, imaju snažno značenje. Ruža je simbol ljubavi, ali i bola. Ona je lepa, ali ima trnje. Isto kao i Erik. Njegova osećanja su iskrena i duboka, ali način na koji ih pokazuje često je pogrešan i zastrašujući. Ruža u loži opere pokazuje da Fantom nije bezosećajno biće, već neko ko voli, sanja i pati, iako to čini iz senke.

U romanu „Fantom iz opere“, Erik je prikazan kao tragična figura, čovek koji je od detinjstva bio odbačen zbog svog izgleda. U operi i mjuziklu ta priča dobija još jaču emotivnu dimenziju kroz muziku, jer muzika postaje njegov jedini pravi jezik. Kroz nju on govori ono što nikada nije smeo da kaže naglas. I u knjizi i u operi, Fantom ostaje neko koga više posmatramo nego što zaista upoznajemo.

Kroz muziku Fantom govori ono što nikada nije smeo da izgovori rečima. Njegove kompozicije nisu samo deo predstave, već odraz njegove unutrašnje borbe, bola i čežnje za ljubavlju koju nije imao. Upravo zato publika često počinje da oseća sažaljenje prema njemu, čak i kada su njegova dela pogrešna i zlonamerna. I u knjizi i u operi, Fantom ostaje lik koga više posmatramo nego što zaista upoznajemo, jer nikada ne ulazimo potpuno u njegov svet. On je stalno na granici – između čoveka i mita, umetnika i čudovišta, Erika i Fantoma.

Zbog toga se postavlja pitanje – da li smo ikada imali priliku da upoznamo pravog Erika, ili smo upoznali samo lik koji je morao da stvori kako bi preživeo? Fantom nije samo uloga koju on igra, već štit koji ga čuva od sveta koji ga je povredio i učinio da se on oseća loše. Njegova priča ostaje snažna upravo zato što nas tera da razmislimo koliko često i sami sudimo ljude po onome što vidimo, a ne po onome što se krije ispod površine. Erik možda nikada ne bi ni postao Fantom da nije doživeo mržnju zbog njegovog izgleda. Možda da je bio prihvaćen, ne bi morao da se krije iza maske, niti da živi ispod opere, daleko od ljudi i svetla pozornice. Možda bi tada Erik bio samo umetnik, kompozitor i čovek koji slobodno pokazuje svoju ljubav i emocije. Umesto toga, društvo ga je nateralo da izabere tamu kao jedino mesto gde može da postoji bez osude.

\end{document}
